\documentclass[12pt]{article}
\usepackage[a4paper,margin=1in]{geometry}
\usepackage[T1]{fontenc}
\usepackage{lmodern,microtype}
\usepackage{amsmath,amssymb,amsthm,mathtools}
\usepackage{booktabs,array,enumitem,xcolor}
\usepackage[numbers,sort&compress]{natbib}
\usepackage[colorlinks=true,linkcolor=blue!55!black,citecolor=blue!55!black,urlcolor=blue!55!black]{hyperref}
\linespread{1.04}

\newtheorem{theorem}{Theorem}[section]
\newtheorem{proposition}[theorem]{Proposition}
\newtheorem{corollary}[theorem]{Corollary}
\theoremstyle{remark}
\newtheorem{remark}[theorem]{Remark}
\newcommand{\Tr}{\operatorname{Tr}}
\newcommand{\cZ}{\mathcal Z}
\newcommand{\norm}[1]{\left\lVert#1\right\rVert}

\title{The Double-Cycle Geometric Cloud\\
\large An Exact Finite Operator Model for the RH-80 Moving Factor}
\author{Liang Wang}
\date{July 2026}

\begin{document}
\maketitle

\begin{abstract}
RH-176 identifies the reduced determinant of a length-$L$ cycle with the
geometric section $\Pi_{L-1}$.  The moving factor required by RH-80 is the
square
\[
 C_N(w)=\Pi_N(w/\lambda)^2,
 \qquad \lambda=1.678573510428322\ldots.
\]
We give an exact finite normal operator realizing this factor.  Let
$L=N+1$, restrict the cycle $C_L$ to the zero-mean subspace, scale by
$\lambda^{-1}$, and take two copies:
\[
 K_N=\lambda^{-1}
 \bigl(C_L^\circ\oplus C_L^\circ\bigr).
\]
Then
\[
 \det(I-wK_N)=\Pi_N(w/\lambda)^2=C_N(w).
\]
The construction has exactly $2N$ nonstationary modes, matching the cloud
count.  Its complete periodic trace ledger is
\[
 \Tr K_N^m=
 \frac{2}{\lambda^m}
 \begin{cases}
 N,&N+1\mid m,\\
 -1,&N+1\nmid m.
 \end{cases}
\]
Consequently the RH-80 geometric factor possesses a strict intrinsic
finite-orbit representation, not merely a formal polynomial identity.

We also recover the local scattering profile
$(N+1)^{-1}\Pi_N(e^{s/(N+1)})\to(e^s-1)/s$.  A 192-case determinant audit and
104 trace checks have maximum errors $9.48\times10^{-15}$ and below
$10^{-10}$, respectively.  This closes the algebraic model factor only.  It
does not identify the actual noisy cloud with $K_N$, select $N(\sigma)$,
construct a physical Riesz projection, or prove complement convergence.
\end{abstract}

\section{The algebraic gap left by RH-80}

RH-80 proves that fixed limiting pole cancellation cannot cross the pole
circle.  Its correct sufficient replacement is division by the actual finite
moving cloud polynomial \citep{WangRH80}.  In the canonical model that factor
is
\begin{equation}\label{eq:cloud}
 C_N(w)=\Pi_N(w/\lambda)^2,
 \qquad
 \Pi_N(q)=1+q+\cdots+q^N.
\end{equation}
The theorem in RH-80 treats $C_N$ as a finite determinant once a reducing
cloud exists, but it does not specify a canonical finite operator whose
determinant is \eqref{eq:cloud}.

RH-176 supplies one copy: a reduced cycle of length $N+1$ has determinant
$\Pi_N$ \citep{WangRH176}.  Two copies therefore appear inevitable.  The
point of this paper is to make that observation operator-theoretic and to
record the resulting traces, multiplicities, and limitations.

\section{Definition of the double-cycle cloud}

Let $L=N+1$, let $C_L$ be the forward cycle on $\mathbb C^L$, and let
$\cZ_L=\mathbf1_L^\perp$.  Write
\[
 C_L^\circ=C_L|_{\cZ_L}.
\]
Its eigenvalues are the nontrivial roots
\[
 \omega_k=e^{2\pi ik/L},
 \qquad 1\le k\le N.
\]
Define the normal $2N$-dimensional operator
\begin{equation}\label{eq:K}
 K_N=\lambda^{-1}
 \bigl(C_L^\circ\oplus C_L^\circ\bigr).
\end{equation}

\begin{theorem}[Exact double-cycle realization]\label{thm:det}
For every $w\in\mathbb C$,
\begin{equation}\label{eq:det}
 \det(I-wK_N)
 =\prod_{k=1}^N(1-w\omega_k/\lambda)^2
 =\Pi_N(w/\lambda)^2.
\end{equation}
Thus $K_N$ realizes the exact RH-80 geometric moving factor.
\end{theorem}

\begin{proof}
Apply the reduced-cycle determinant identity to each direct-sum copy:
\[
 \det(I-w\lambda^{-1}C_L^\circ)
 =\frac{1-(w/\lambda)^L}{1-w/\lambda}
 =\Pi_N(w/\lambda).
\]
Determinants multiply over finite invariant direct sums, proving
\eqref{eq:det}.
\end{proof}

\begin{corollary}[Exact dimension dictionary]
The model has cycle length $L=N+1$, reduced rank $N$ per copy, total cloud
rank $2N$, and determinant degree $2N$.
\end{corollary}

This agrees with the RH-15 convention that an effective degree $N$ selects
$2N$ parity-resolved cloud eigenvalues \citep{WangRH15}.  Agreement of counts
does not identify the observed eigenvalues with the exact roots.

\section{Complete periodic trace ledger}

The cycle powers have
\[
 \Tr C_L^m=
 \begin{cases}
 L,&L\mid m,\\
 0,&L\nmid m.
 \end{cases}
\]
Removing the stationary mode subtracts one from every power trace.

\begin{theorem}[Double-cycle trace formula]\label{thm:trace}
For every integer $m\ge1$,
\begin{equation}\label{eq:trace}
 \Tr K_N^m=
 \frac{2}{\lambda^m}
 \begin{cases}
 N,&N+1\mid m,\\
 -1,&N+1\nmid m.
 \end{cases}
\end{equation}
Equivalently,
\begin{equation}\label{eq:logdet}
 \log\det(I-wK_N)
 =-\sum_{m\ge1}\frac{w^m}{m}\Tr K_N^m
\end{equation}
for $|w|<\lambda$.
\end{theorem}

\begin{proof}
On one reduced cycle,
\[
 \Tr(C_L^\circ)^m=\Tr C_L^m-1,
\]
which is $N$ when $L\mid m$ and $-1$ otherwise.  Scaling contributes
$\lambda^{-m}$ and doubling contributes the factor two.  The finite
log-determinant expansion is standard inside the reciprocal spectral radius
\citep{Simon2005}.
\end{proof}

The negative background value $-2\lambda^{-m}$ is not an error.  It is the
trace contribution of deleting two stationary modes.  At multiples of the
cycle length, all phases align and the trace jumps to
$2N\lambda^{-m}$.

\section{Coefficient ledger of the squared geometric section}

Writing
\begin{equation}\label{eq:coefficients}
 \Pi_N(q)^2=\sum_{j=0}^{2N}c_{N,j}q^j,
\end{equation}
the coefficients count pairs $(a,b)$ with
$0\le a,b\le N$ and $a+b=j$.

\begin{proposition}[Triangular coefficient law]\label{prop:coefficients}
The exact coefficients are
\begin{equation}\label{eq:triangular}
 c_{N,j}=
 \begin{cases}
 j+1,&0\le j\le N,\\
 2N-j+1,&N<j\le2N.
 \end{cases}
\end{equation}
Therefore
\begin{equation}\label{eq:scaled-coefficients}
 \det(I-wK_N)
 =\sum_{j=0}^{2N}c_{N,j}\lambda^{-j}w^j.
\end{equation}
\end{proposition}

\begin{proof}
For $j\le N$, the admissible pairs are
$(0,j),(1,j-1),\ldots,(j,0)$, giving $j+1$.  For $j>N$, the square
$[0,N]^2$ truncates both ends and leaves $2N-j+1$ pairs.  Substitute
$q=w/\lambda$.
\end{proof}

The coefficient law gives a finite target for the physical Q interface.  A
future Riesz cloud need not match individual exact roots before one can test
its determinant: its normalized coefficients can first be compared with the
triangular ledger \eqref{eq:scaled-coefficients}.  Newton identities make this
equivalent to matching the power traces in Theorem~\ref{thm:trace}, but the
two diagnostics have different numerical conditioning.

\section{One-step square roots and the two-step double cycle}

RH-15 uses the one-step cloud
\begin{equation}\label{eq:one-step-roots}
 \mu_{N,k}^{\pm}=\lambda^{-1/2}
 \exp\left(\pm\frac{ik\pi}{N+1}\right),
 \qquad 1\le k\le N.
\end{equation}
There are $2N$ roots and their regularized determinant is
$\Pi_N(z^2/\lambda)$ \citep{WangRH15}.

\begin{theorem}[Squaring produces the double-cycle multiset]
\label{thm:squaring}
The multiset of squares of the one-step roots
\eqref{eq:one-step-roots} is exactly the spectrum of $K_N$:
\begin{equation}\label{eq:squared-roots}
 \{(\mu_{N,k}^{\pm})^2:1\le k\le N\}
 =\biguplus_{j=1}^N
 \left\{\lambda^{-1}e^{2\pi ij/(N+1)},
 \lambda^{-1}e^{2\pi ij/(N+1)}\right\}
\end{equation}
with every nontrivial root appearing twice.
\end{theorem}

\begin{proof}
Squaring gives
$\lambda^{-1}e^{\pm2\pi ik/(N+1)}$.  As $k$ runs from one to $N$, the plus
sign lists all nontrivial roots once, and the minus sign lists the same set in
reverse order.  Hence each root has multiplicity two.
\end{proof}

This theorem supplies an exact algebraic adapter between the RH-15 one-step
cloud and the RH-80 two-step squared factor.  It is consistent with the
general identity relating a one-step regularized determinant to its symmetric
two-step Fredholm completion.  It still does not prove that the physical
one-step resonances square without leakage into the selected physical
two-step Riesz cloud; that is an operator-selection problem, not a polynomial
identity.

\section{Normality, uniqueness, and orientation blindness}

The operator $K_N$ is normal.  Among finite normal operators, its determinant
fixes the multiset of eigenvalues and therefore fixes it up to unitary
equivalence.  This is a useful limited canonicity: once $N$, $\lambda$, and
the squared geometric factor are specified, no further normal spectral data
are missing.

However, orientation is not spectral data.  Replacing either cycle by its
inverse conjugates and permutes the root set, leaving
\eqref{eq:det} and all ordinary traces unchanged.  The same scalar cloud
factor can therefore arise from
\[
 C_L^\circ\oplus C_L^\circ,
 \qquad
 C_L^\circ\oplus(C_L^\circ)^{-1},
 \qquad\text{or}\qquad
 (C_L^\circ)^{-1}\oplus(C_L^\circ)^{-1}.
\]
These completions carry different time arrows but the same determinant.
RH-178 adds a rank-one directed marker to separate them.

\section{Scattering coordinate}

The finite section has a nontrivial edge scaling.  Put
$q=e^{s/(N+1)}$.  Then
\begin{equation}\label{eq:profile}
 \frac1{N+1}\Pi_N(q)
 =\frac{e^s-1}{(N+1)(e^{s/(N+1)}-1)}.
\end{equation}

\begin{proposition}[Geometric scattering profile]\label{prop:profile}
For every $s\in\mathbb C$,
\begin{equation}\label{eq:limit}
 \frac1{N+1}\Pi_N(e^{s/(N+1)})
 \longrightarrow \frac{e^s-1}{s},
\end{equation}
with the removable value one at $s=0$, locally uniformly in $s$.
The normalized double-cycle determinant profile is the square of this limit.
\end{proposition}

\begin{proof}
The denominator in \eqref{eq:profile} converges locally uniformly to $s$.
The removable value follows by differentiation or continuity.  Squaring
preserves locally uniform convergence on compact sets.
\end{proof}

This is the same canonical profile used in the RH-15 scattering analysis.
Here it arises from a literal finite normal operator, not only from a
polynomial approximation.

\section{Audit}

The audit uses degrees
\[
 N=1,2,3,4,7,12,20,32
\]
and 24 random complex spectral parameters per degree inside
$|w|<0.8\lambda$, giving 192 determinant cases.  The maximum relative error
between the eigenvalue product and $\Pi_N(w/\lambda)^2$ is
$9.47\times10^{-15}$.

For each degree it also checks powers up to
$\min(2(N+1),20)$, for 104 trace cases.  The maximum absolute difference from
\eqref{eq:trace} is below $10^{-10}$.  Finally, the five edge coordinates
$s=-1,-1/2,0,1/2,1$ are evaluated; at $N=32$ the finite profile is already
close to $(e^s-1)/s$.  These are implementation checks of exact finite
identities.

\section{Relation to the physical Gate-A program}

The theorem closes an algebraic leaf that may be denoted
$Q_{\rm geom}$:
\[
 \text{squared geometric cloud factor}
 \Longleftrightarrow
 \text{double reduced cycle determinant}.
\]
It does not close the physical cloud ledger $Q_{\rm phys}$.  The latter must
prove that the actual noisy Riesz cloud has the same divisor or converges to
it in the required coefficient sense.  Nor does this paper provide the
physical packet-to-Riesz map $R$, the complementary Schatten limit $U$, the
normalization leaf $Z$, or the directed limit $T$.

The next two concrete tasks are:
\begin{enumerate}[leftmargin=2em]
\item supplement the scalar determinant with orientation-separating marked
traces;
\item compare $L=N+1$ with the reset clock rank and observed cloud degree.
\end{enumerate}
No self-adjoint Hilbert--Polya operator, prime-power trace formula, zeta
divisor identity, or Riemann Hypothesis is claimed.

\bibliographystyle{plainnat}
\bibliography{references}
\end{document}
